\documentclass[11pt]{article}
\usepackage{mathblatt}
\usepackage[hidelinks]{hyperref}
\newcommand{\zweigzeile}[1]{\par\nopagebreak{\small #1}\par\medskip}
\begin{document}
\blattkopf{Lineare Funktionen}{Lernblatt}
\weit
\hypertarget{e1}{}\einheitenkopf{Einheit 1 von 5 · Proportionale Funktion}
\zweigzeile{Hier lernst du etwas · neu in diesem Jahr · P10 oft · baut auf: Dreisatz}
\begin{aufgabe}{\textbf{Ich kann ein Steigungsdreieck lesen.} Test.}
\begin{ksys}[xmin=-4,xmax=4,ymin=-4,ymax=4,ablesen]
\gerade{2}{-1}{g}\steigungsdreieck{0}{-1}{2}
\gerade[-2]{-0.5}{1}{h}\steigungsdreieck{0}{1}{-0.5}
\end{ksys}

\begin{teile}
\teil Test $f(x) = 2x$ \feld{m}\feld{n}\feldl{y}
\teil \kreuz{steigt}\kreuz{fällt}
\end{teile}
\punktprobenliste{2x-1/A(3|5), -\frac{3}{4}x+2/B(-8|8)}
\mnliste{2x+3, 0{,}5x-1}
\wertetabelle{x}{f(x)}{-2,-1,0,1,2}
\begin{geruest} \gz{a}{Gerade $g$}{\feld{m}\feld{n}\feldl{y}} \end{geruest}
\end{aufgabe}
\begin{aufgabe}{Zweite}
\beispiel{3x + 2 &= 14 &&\mid -2 \\ 3x &= 12 &&\mid :3 \\ x &= 4}
\begin{gleichungsraster}[2] \gl{2x + 4 = 10} & \gl{3x - 5 = 7} \\ \end{gleichungsraster}
\begin{dreisatz}{$x$}{$y$}
\dsz{2}{3}
\dsp{\cdot 6}{\cdot 6}
\dsz{12}{\dsleer}
\end{dreisatz}
\end{aufgabe}
\newpage
\begleitteil
\erg{1}{a) test \quad b) test}
\end{document}
